\chapterimage{chapter_head_1.pdf}
\chapter{C Review \& Dynamic Memory Allocation}

\section{Theoretical Core}

\begin{definition}[Pointers]
A variable that stores a memory address as its value. Pointers allow direct manipulation of memory, efficient data sharing between functions (pass by reference), and referring to large data structures compactly. The \texttt{\&} operator (address-of) retrieves the memory address of a variable, and the \texttt{*} operator (dereference) accesses or manipulates the value stored at that location.
\end{definition}

\begin{definition}[Structs]
A custom data type defined using the \texttt{struct} keyword that allows grouping related variables (members/fields) of different data types into a single unit. It is often used alongside \texttt{typedef} to create aliases for easier reference.
\end{definition}

\begin{definition}[Dynamic Memory Allocation (DMA)]
Allocating memory at runtime from the heap space when the required size is unknown at compile time.
\begin{itemize}
    \item \texttt{malloc()}: Allocates a specified number of bytes on the heap and returns a void pointer to it. The allocated memory is uninitialized.
    \item \texttt{calloc()}: Allocates memory for an array of elements and initializes all bytes to zero.
    \item \texttt{realloc()}: Resizes a previously allocated memory block, preserving its content up to the minimum of the old and new sizes.
    \item \texttt{free()}: Deallocates dynamically allocated memory, releasing it back to the heap to prevent memory leaks.
\end{itemize}
\end{definition}

\newpage
\section{Algorithmic Tracing}

\begin{example}[Dynamically Allocating a 2D Array]
\begin{enumerate}
    \item Determine the number of rows and columns needed.
    \item Define a double pointer variable (e.g., \texttt{int **arr}).
    \item Call \texttt{malloc()} to allocate an array of pointers (one for each row) and assign the result to the double pointer: \texttt{arr = malloc(rows * sizeof(int *));}.
    \item Iterate over each row pointer (\texttt{i} from 0 to \texttt{rows - 1}).
    \item Inside the loop, call \texttt{malloc()} to allocate memory for the columns and assign it to the current row pointer: \texttt{arr[i] = malloc(cols * sizeof(int));}.
\end{enumerate}
\end{example}

\begin{example}[Freeing a 2D Array]
\begin{enumerate}
    \item Iterate over each row pointer (\texttt{i} from 0 to \texttt{rows - 1}).
    \item Call \texttt{free()} on each row array: \texttt{free(arr[i]);}.
    \item Finally, call \texttt{free()} on the array of pointers itself: \texttt{free(arr);}.
\end{enumerate}
\end{example}

\newpage
\section{Code Implementation}
The following snippet demonstrates dynamic memory allocation (including a 2D array) and the correct usage of struct member access operators.

\begin{lstlisting}[language=C]
// Struct definitions demonstrating the dot (.) and arrow (->) operators
typedef struct Song_s {
    int song_id;
    char title[101];
} Song;

typedef struct Cat_s {
    char *name;         // Dynamically allocated string
    int age;
} Cat;

void allocate_and_access() {
    // Dot operator (.) usage with a struct instance
    Song s;
    s.song_id = 101;

    // Arrow operator (->) usage with a pointer to a struct
    Cat *cat = malloc(sizeof(Cat));
    cat->age = 5;
    
    // Remember to free dynamically allocated struct instances
    free(cat);
}

// 2D Array Dynamic Allocation & Deallocation
void manage_2d_array(int rows, int cols) {
    // 1. Allocation
    int **matrix = malloc(rows * sizeof(int *));
    for (int i = 0; i < rows; i++) {
        matrix[i] = malloc(cols * sizeof(int));
    }
    
    // 2. Deallocation
    for (int i = 0; i < rows; i++) {
        free(matrix[i]); // Free each row
    }
    free(matrix); // Free the array of pointers
}
\end{lstlisting}

\section{Skills \& Pitfalls}

\subsection{Must-Know Skills}
\begin{itemize}
    \item \textbf{Arrow vs. Dot Operator}: Use the dot operator (\texttt{.}) to access members of a struct variable (e.g., \texttt{list.count}). Use the arrow operator (\texttt{->}) to access members when you have a pointer to a struct (e.g., \texttt{list->count}), which is shorthand for \texttt{(*list).count}.
    \item \textbf{Pointer Arithmetic}: Understand the precedence of pointer arithmetic. \texttt{*(arr + i)} correctly dereferences the \texttt{i}-th element, whereas \texttt{*arr + i} dereferences the first element and adds \texttt{i} to its value.
    \item \textbf{Constructors and Destructors}: Modularize code by defining helper functions to handle the allocation (e.g., \texttt{createStore}) and deallocation (e.g., \texttt{freeStore}) of complex dynamic structures to reduce the risk of memory leaks.
\end{itemize}

\subsection{Common Mistakes (Pitfalls)}
\begin{remark}
\begin{itemize}
    \item \textbf{Memory Leaks}: Occur when dynamically allocated memory is no longer needed but is not released using \texttt{free()}. Every \texttt{malloc()} should have a corresponding \texttt{free()}.
    \item \textbf{Dangling Pointers}: Continuing to use a pointer after the memory it points to has been freed.
    \item \textbf{Segmentation Faults}: Run-time errors caused by memory access violations, such as dereferencing an invalid memory location or an uninitialized pointer. Always initialize pointers to \texttt{NULL} and check them defensively (e.g., \texttt{if (ptr != NULL)}) before dereferencing.
\end{itemize}
\end{remark}
